\section{Chapter 1: Basic Concepts in Astrophysics}

\subsection{Questions}

\begin{enumerate}[label=1.\arabic*]
    \item The star Gliese $710$ is presently $62$ lightyears from Earth, but Hipparcos data suggest that in about a million years it will pass within one lightyear of Earth. Its appearent magnitude is $9.7$.
    \begin{enumerate}[label=(\alph*)]
        \item What is the absolute magnitude of Gliese $710$?
        \item What will be its apparent visual magnitude be in a million years as viewed from Earth?  
    \end{enumerate}   
    \item Estimate the hydrodynamical free fall timescale for the Sun, a red giant of $M=M_{\odot}$ and $R=100R_{\odot}$, and a white dwarf of $M=M_{\odot}$ and $R=0.01R_{\odot}$ and $R=100R_{\odot}$, 
    and a white dwarf of $M=M_{\odot}$ and $R=0.02R_{\odot}$. 
    \item Suppose that a star of mass $M$ and radius $R$ has a density distribution $\rho(r) = \rho_c(1 - \frac{r}{R})$,
where $\rho_c$ is the density at the center of the star.
    \begin{enumerate}[label=(\alph*)]
        \item Find an expression for $\rho_c$ in terms of $M$ and $R$.
        \item Calculate the mass $m(r)$ interior to radius $r$.
        \item Calculate the total gravitational binding energy of the star.
        \item Find an expression for the pressure $P(r)$ as a function of $r$, assuming that the star is in hydrostatic equilibrium and that the pressure at the surface $(r = R)$ is zero.
        \item Assume that the material in the star is a \textbf{monatomic} ideal gas. Calculate the total
                internal energy of the star from $P(r)$, and show that the virial theorem is satisfied.
    \end{enumerate}
\end{enumerate}

\subsection{Solutions}
\subsubsection*{Solution 1.1: Absolute magnitude and apparent magnitude of Gliese $710$}
\begin{enumerate}[label=(\alph*)]
    \item The absolute magnitude $M$ of a star can be calculated using the formula:
    \begin{equation}
    M = m - 5 \log_{10} \left( \frac{d}{10} \right),
    \end{equation}
    where $m$ is the apparent magnitude and $d$ is the distance to the star in parsecs. Since $1$ lightyear is approximately $0.3064$ parsecs, we can convert the distance to parsecs:
    \begin{equation}
    d = 62 \text{ lightyears} \times 0.3064 \text{ parsecs/lightyear} \approx 19.01 \text{ parsecs}.
    \end{equation}
    Now we can calculate the absolute magnitude:
    \begin{align*}
    M &= 9.7 - 5 \log_{10} \left( \frac{19.01}{10} \right) \\
      &= 9.7 - 5 \log_{10} (1.901) \\
      &\approx 9.7 - 5 \times 0.279 = 9.7 - 1.395 = 8.305.
    \end{align*}
    Thus, the absolute magnitude of Gliese $710$ is approximately $8.3$.
    
    \item The apparent magnitude $m'$ of Gliese $710$ in a million years can be calculated using the same formula, but with the new distance of $1$ lightyear (or $0.3064$ parsecs):
    \begin{align*}
    m' &= M + 5 \log_{10} \left( \frac{d'}{10} \right) \\
       &= 8.305 + 5 \log_{10} (0.3064) \\
       &\approx 8.305 + 5 (-0.513) = 8.305 - 2.565 = 5.74.
    \end{align*}
    Thus, the apparent visual magnitude of Gliese $710$ in a million years will be approximately $5.7$.
\end{enumerate}

\subsubsection*{Solution 1.2: Hydrodynamical free fall timescale}
The hydrodynamical free fall timescale $t_{ff}$ can be estimated using the formula:
\begin{equation}
t_{ff} = \sqrt{\frac{3\pi}{32 G \rho}},
\end{equation}
where $G$ is the gravitational constant and $\rho$ is the average density of the star. The average density can be calculated using the formula:
\begin{equation}
\rho = \frac{3M}{4\pi R^3}.
\end{equation}  
For the Sun, we have $M = M_\odot$ and $R = R_\odot$, so the average density is:
\begin{equation}
\rho_\odot = \frac{3M_\odot}{4\pi R_\odot^3}.
\end{equation}
Substituting this into the formula for $t_{ff}$, we get:
\begin{equation}
t_{ff,\odot} = \sqrt{\frac{3\pi}{32 G \rho_\odot}} = \sqrt{\frac{3\pi}{32 G \frac{3M_\odot}{4\pi R_\odot^3}}} = \sqrt{\frac{4\pi^2 R_\odot^3}{32 G M_\odot}} = \sqrt{\frac{\pi^2 R_\odot^3}{8 G M_\odot}}.
\end{equation}
For a red giant with $M = M_\odot$ and $R = 100 R_\odot$, the average density is:
\begin{equation}
\rho_{RG} = \frac{3M_\odot}{4\pi (100 R_\odot)^3} = \frac{3M_\odot}{4\pi 10^6 R_\odot^3} = \frac{3M_\odot}{4\pi 10^6 R_\odot^3}.
\end{equation}
Substituting this into the formula for $t_{ff}$, we get:
\begin{equation}
t_{ff,RG} = \sqrt{\frac{3\pi}{32 G \rho_{RG}}} = \sqrt{\frac{3\pi}{32 G \frac{3M_\odot}{4\pi 10^6 R_\odot^3}}} = \sqrt{\frac{4\pi^2 10^6 R_\odot^3}{32 G M_\odot}} = \sqrt{\frac{\pi^2 10^6 R_\odot^3}{8 G M_\odot}}.
\end{equation}
For a white dwarf with $M = M_\odot$ and $R = 0.02 R_\odot$, the average density is:
\begin{equation}
\rho_{WD} = \frac{3M_\odot}{4\pi (0.02 R_\odot)^3} = \frac{3M_\odot}{4\pi 8 \times 10^{-6} R_\odot^3} = \frac{3M_\odot}{4\pi 8 \times 10^{-6} R_\odot^3}.
\end{equation}
Substituting this into the formula for $t_{ff}$, we get:
\begin{equation}
t_{ff,WD} = \sqrt{\frac{3\pi}{32 G \rho_{WD}}} = \sqrt{\frac{3\pi}{32 G \frac{3M_\odot}{4\pi 8 \times 10^{-6} R_\odot^3}}} = \sqrt{\frac{4\pi^2 8 \times 10^{-6} R_\odot^3}{32 G M_\odot}} = \sqrt{\frac{\pi^2 8 \times 10^{-6} R_\odot^3}{8 G M_\odot}}.
\end{equation}
Thus, the hydrodynamical free fall timescales for the Sun, the red giant, and the white dwarf can be calculated using the above formulas.

\subsubsection*{Solution 1.3: Density distribution and hydrostatic equilibrium}
\begin{enumerate}[label=(\alph*)]
    \item To find an expression for $\rho_c$ in terms of $M$ and $R$, we can integrate the density distribution over the volume of the star to find the total mass:
    \begin{equation}
    M = \int_0^R 4\pi r^2 \rho(r) dr = \int_0^R 4\pi r^2 \rho_c \left(1 - \frac{r}{R}\right) dr.
    \end{equation}
    Evaluating this integral, we get:
    \begin{align*}
    M &= 4\pi \rho_c \int_0^R r^2 \left(1 - \frac{r}{R}\right) dr \\
      &= 4\pi \rho_c \left[ \int_0^R r^2 dr - \frac{1}{R} \int_0^R r^3 dr \right] \\
      &= 4\pi \rho_c \left[ \frac{R^3}{3} - \frac{1}{R} \cdot \frac{R^4}{4} \right] = 4\pi \rho_c R^3 \left( \frac{1}{3} - \frac{1}{4} \right) = 4\pi \rho_c R^3 \cdot \frac{1}{12} = \frac{\pi}{3} \rho_c R^3.
    \end{align*}
    Solving for $\rho_c$, we get:
    \begin{equation}
    \rho_c = \frac{3M}{\pi R^3}.
    \end{equation}
    \item To calculate the mass $m(r)$ interior to radius $r$, we can integrate the density distribution from $0$ to $r$:
    \begin{equation}
    m(r) = \int_0^r 4\pi r'^2 \rho(r') dr' = \int_0^r 4\pi r'^2 \rho_c \left(1 - \frac{r'}{R}\right) dr'.
    \end{equation}
    Evaluating this integral, we get:
    \begin{align*}
    m(r) &= 4\pi \rho_c \int_0^r r'^2 \left(1 - \frac{r'}{R}\right) dr' \\
         &= 4\pi \rho_c \left[ \int_0^r r'^2 dr' - \frac{1}{R} \int_0^r r'^3 dr' \right] \\
         &= 4\pi \rho_c \left[ \frac{r^3}{3} - \frac{1}{R} \cdot \frac{r^4}{4} \right] = 4\pi \rho_c r^3 \left( \frac{1}{3} - \frac{r}{4R} \right). 
    \end{align*}
    Substituting $\rho_c$ from part (a), we get:
    \begin{equation}
    m(r) = 4\pi \cdot \frac{3M}{\pi R^3} r^3 \left( \frac{1}{3} - \frac{r}{4R} \right) = \frac{12M}{R^3} r^3 \left( \frac{1}{3} - \frac{r}{4R} \right).
    \end{equation}
    \item To calculate the total gravitational binding energy of the star, we can use the formula:
    \begin{equation}
        U = -\int_0^R \frac{G m(r)}{r} 4\pi r^2 \rho(r) dr.
    \end{equation}
    Substituting $m(r)$ and $\rho(r)$, we get:
    \begin{align*}
        U &= -\int_0^R \frac{G}{r} \cdot \frac{12M}{R^3} r^3 \left( \frac{1}{3} - \frac{r}{4R} \right) \cdot \rho_c \left(1 - \frac{r}{R}\right) dr \\
          &= -\int_0^R \frac{G}{r} \cdot \frac{12M}{R^3} r^3 \left( \frac{1}{3} - \frac{r}{4R} \right) \cdot \frac{3M}{\pi R^3} \left(1 - \frac{r}{R}\right) dr \\
          &= -\frac{36G M^2}{\pi R^6} \int_0^R r^2 \left( \frac{1}{3} - \frac{r}{4R} \right) \left(1 - \frac{r}{R}\right) dr.
    \end{align*}
    Evaluating this integral will give the total gravitational binding energy of the star.
    \item To find an expression for the pressure $P(r)$ as a function of $r$, we can use the equation of hydrostatic equilibrium:
    \begin{equation}
    \frac{dP}{dr} = -\frac{G m(r) \rho(r)}{r^2}.
    \end{equation}
    Substituting $m(r)$ and $\rho(r)$, we get:
    \begin{align*}
        \frac{dP}{dr} &= -\frac{G}{r^2} \cdot \frac{12M}{R^3} r^3 \left( \frac{1}{3} - \frac{r}{4R} \right) \cdot \rho_c \left(1 - \frac{r}{R}\right) \\
                       &= -\frac{12G M}{R^3} r \left( \frac{1}{3} - \frac{r}{4R} \right) \cdot \frac{3M}{\pi R^3} \left(1 - \frac{r}{R}\right) \\
                       &= -\frac{36G M^2}{\pi R^6} r \left( \frac{1}{3} - \frac{r}{4R} \right) \left(1 - \frac{r}{R}\right).
    \end{align*}
    Integrating this equation will give the pressure as a function of $r$.
    \item Assuming that the material in the star is a monatomic ideal gas, the internal energy $U$ can be calculated using the formula:
    \begin{equation}    
        U = \frac{3}{2} \int_0^R 4\pi r^2 P(r) dr.
    \end{equation}
    Substituting $P(r)$ from part (d), we can evaluate this integral to find the total internal energy of the star. To show that the virial theorem is satisfied, we need to verify that $2U + U_{grav} = 0$, where $U_{grav}$ is the gravitational binding energy calculated in part (c).
\end{enumerate}